% Part: first-order-logic
% Chapter: syntax-and-semantics
% Section: Semantic Notions

\documentclass[../../../include/open-logic-section]{subfiles}

\begin{document}

\olfileid{fol}{syn}{sem}
\olsection{અર્થવિચારી ખ્યાલો}

\begin{explain}
પ્રથમ-ક્રમ ભાષાઓ માટેની !!{structure}sની વ્યાખ્યા આપ્યા પછી, આપણે
વાક્યોના કેટલાક પાયાના અર્થવિચારી ગુણધર્મો અને વાક્યો વચ્ચેના સંબંધો
વ્યાખ્યાયિત કરી શકીએ છીએ. એમાં સૌથી સરળ ખ્યાલ વાક્યના
\emph{પ્રામાણ્ય}નો છે. કોઈ વાક્ય દરેક !!{structure}માં સંતોષાય તો તે
પ્રમાણભૂત છે. પ્રમાણભૂત વાક્યો એવાં વાક્યો છે જે તેમાંના અતાર્કિક
સંકેતોનું ગમે તે રીતે અર્થઘટન કરીએ તોપણ સંતોષાય છે. તેથી પ્રમાણભૂત
વાક્યોને \emph{તાર્કિક સત્યો} પણ કહે છે---તે કોઈ પણ !!{structure}માં
સત્ય, એટલે કે સંતોષાયેલાં, હોય છે; તેથી તેમની સત્યતા માત્ર તેમાં આવતા
તાર્કિક સંકેતો અને તેમની વાક્યરચનાત્મક !!{structure} પર આધાર રાખે છે,
અતાર્કિક સંકેતો કે તેમના અર્થઘટન પર નહિ.
\end{explain}

\begin{defn}[પ્રામાણ્ય]
વાક્ય~$!A$ \emph{પ્રમાણભૂત}, $\Entails !A$, ત્યારે અને માત્ર ત્યારે
જ્યારે દરેક !!{structure}~$\Struct M$ માટે $\Sat{M}{!A}$.
\end{defn}

\begin{defn}[ફલિતતા]
વાક્યોનો ગણ~$\Gamma$, વાક્ય~$!A$ને \emph{ફલિત કરે} છે,
$\Gamma\Entails !A$, ત્યારે અને માત્ર ત્યારે જ્યારે
$\Sat{M}{\Gamma}$ હોય એવી દરેક !!{structure}~$\Struct M$ માટે
$\Sat{M}{!A}$.
\end{defn}


\begin{defn}[સંતોષ્યતા]
વાક્યોનો ગણ~$\Gamma$ \emph{સંતોષ્ય} છે જો કોઈ !!{structure}
~$\Struct M$ માટે $\Sat{M}{\Gamma}$. જો~$\Gamma$ સંતોષ્ય ન હોય, તો તેને
\emph{અસંતોષ્ય} કહે છે.
\end{defn}

\begin{prop}
વાક્ય~$!A$ પ્રમાણભૂત ત્યારે અને માત્ર ત્યારે જ્યારે વાક્યોના દરેક
ગણ~$\Gamma$ માટે $\Gamma\Entails !A$.
\end{prop}

\begin{proof}
આગળની દિશા માટે, $!A$ પ્રમાણભૂત હોય અને~$\Gamma$ વાક્યોનો ગણ હોય.
$\Sat{M}{\Gamma}$ હોય એવી !!a{structure}~$\Struct M$ લો. $!A$
પ્રમાણભૂત હોવાથી $\Sat{M}{!A}$, એટલે $\Gamma\Entails !A$.

ઊલટી દિશાના પ્રતિનિષેધ માટે, $!A$ અપ્રમાણભૂત હોય; એટલે એવી
!!a{structure}~$\Struct M$ છે કે $\Sat/{M}{!A}$. જ્યારે
$\Gamma=\{\ltrue\}$, ત્યારે~$\ltrue$ પ્રમાણભૂત હોવાથી
$\Sat{M}{\Gamma}$. તેથી એવી !!a{structure}~$\Struct M$ છે કે
$\Sat{M}{\Gamma}$ પરંતુ $\Sat/{M}{!A}$; એટલે~$\Gamma$,~$!A$ને ફલિત
કરતું નથી.
\end{proof}

\begin{prop}
\ollabel{prop:entails-unsat}
$\Gamma\Entails !A$ ત્યારે અને માત્ર ત્યારે જ્યારે
$\Gamma\cup\{\lnot !A\}$ અસંતોષ્ય છે.
\end{prop}

\begin{proof}
આગળની દિશા માટે, ધારો કે $\Gamma\Entails !A$ અને, વિરુદ્ધતાની ખાતર,
એવી !!a{structure}~$\Struct M$ છે કે
$\Sat{M}{\Gamma\cup\{\lnot !A\}}$. $\Sat{M}{\Gamma}$ અને
$\Gamma\Entails !A$ હોવાથી $\Sat{M}{!A}$. વળી,
$\Sat{M}{\Gamma\cup\{\lnot !A\}}$ હોવાથી $\Sat{M}{\lnot !A}$; એટલે
આપણને $\Sat{M}{!A}$ અને $\Sat/{M}{!A}$ બંને મળે છે, જે વિસંવાદ છે.
તેથી આવી કોઈ !!{structure}~$\Struct M$ હોઈ શકે નહિ, એટલે
$\Gamma\cup\{\lnot !A\}$ અસંતોષ્ય છે.

ઊલટી દિશા માટે, ધારો કે $\Gamma\cup\{\lnot !A\}$ અસંતોષ્ય છે. તેથી
દરેક !!{structure}~$\Struct M$ માટે કાં તો $\Sat/{M}{\Gamma}$ અથવા
$\Sat{M}{!A}$. પરિણામે, $\Sat{M}{\Gamma}$ હોય એવી દરેક
!!{structure}~$\Struct M$ માટે $\Sat{M}{!A}$; એટલે
$\Gamma\Entails !A$.
\end{proof}

\begin{prob}
\begin{enumerate}
\item $\Gamma\Entails\bot$ ત્યારે અને માત્ર ત્યારે જ્યારે~$\Gamma$
  અસંતોષ્ય છે, તે બતાવો.
\item $\Gamma\cup\{!A\}\Entails\bot$ ત્યારે અને માત્ર ત્યારે જ્યારે
  $\Gamma\Entails\lnot !A$, તે બતાવો.
\item ધારો કે $c$,~$!A$ કે~$\Gamma$માં આવતું નથી. બતાવો કે
  $\Gamma\Entails\lforall[x][!A]$ ત્યારે અને માત્ર ત્યારે જ્યારે
  $\Gamma\Entails\Subst{!A}{c}{x}$.
\end{enumerate}
\end{prob}

\begin{prop}
જો $\Gamma\subseteq\Gamma'$ અને $\Gamma\Entails !A$ હોય, તો
$\Gamma'\Entails !A$.
\end{prop}

\begin{proof}
ધારો કે $\Gamma\subseteq\Gamma'$ અને $\Gamma\Entails !A$.
$\Sat{M}{\Gamma'}$ હોય એવી સંરચના~$\Struct M$ લો; ત્યારે
$\Sat{M}{\Gamma}$, અને $\Gamma\Entails !A$ હોવાથી $\Sat{M}{!A}$.
આથી જ્યારે પણ $\Sat{M}{\Gamma'}$, ત્યારે $\Sat{M}{!A}$; એટલે
$\Gamma'\Entails !A$.
\end{proof}


\begin{thm}[અર્થાનુસારી નિગમન પ્રમેય]
\ollabel{thm:sem-deduction}
$\Gamma\cup\{!A\}\Entails !B$ ત્યારે અને માત્ર ત્યારે જ્યારે
$\Gamma\Entails !A\lif !B$.
\end{thm}

\begin{proof}
આગળની દિશા માટે, $\Gamma\cup\{!A\}\Entails !B$ હોય અને
$\Sat{M}{\Gamma}$ હોય એવી !!a{structure}~$\Struct M$ લો. જો
$\Sat{M}{!A}$, તો $\Sat{M}{\Gamma\cup\{!A\}}$; અને
$\Gamma\cup\{!A\}$,~$!B$ને ફલિત કરતું હોવાથી આપણને
$\Sat{M}{!B}$ મળે છે. તેથી $\Sat{M}{!A\lif !B}$, એટલે
$\Gamma\Entails !A\lif !B$.

ઊલટી દિશા માટે, $\Gamma\Entails !A\lif !B$ હોય અને
$\Sat{M}{\Gamma\cup\{!A\}}$ હોય એવી !!a{structure}~$\Struct M$ લો.
ત્યારે $\Sat{M}{\Gamma}$, એટલે $\Sat{M}{!A\lif !B}$; અને
$\Sat{M}{!A}$ હોવાથી $\Sat{M}{!B}$. આથી જ્યારે પણ
$\Sat{M}{\Gamma\cup\{!A\}}$, ત્યારે $\Sat{M}{!B}$; એટલે
$\Gamma\cup\{!A\}\Entails !B$.
\end{proof}

\begin{prop}\ollabel{prop:quant-terms}
  $\Struct{M}$ !!a{structure}, એક મુક્ત ચલ~$x$ ધરાવતું~$!A(x)$
  !!a{formula} અને~$t$ બંધ પદ હોય. ત્યારે:
  \begin{tagenumerate}{prvEx,prvAll}
    \tagitem{prvEx}{$!A(t)\Entails\lexists[x][!A(x)]$}{}
    \tagitem{prvAll}{$\lforall[x][!A(x)]\Entails !A(t)$}{}
  \end{tagenumerate}
\end{prop}

\begin{proof}
  \begin{tagenumerate}{prvEx,prvAll}
    \tagitem{prvEx}{%
      \iftag{probEx}{મહાવરો.}{ધારો કે $\Sat{M}{!A(t)}$. એવી
        ચલ-નિયુક્તિ~$s$ લો કે $s(x)=\Value{t}{M}$. $!A(t)$
        !!a{sentence} હોવાથી $\Sat{M}{!A(t)}[s]$. હવે
        \olref[ext]{prop:ext-formulas}થી $\Sat{M}{!A(x)}[s]$.
        \olref[ass]{prop:sat-quant}થી
        $\Sat{M}{\lexists[x][!A(x)]}$.}}{}
    \tagitem{prvAll}{%
      \iftag{probAll}{મહાવરો.}{ધારો કે
        $\Sat{M}{\lforall[x][!A(x)]}$. એવી ચલ-નિયુક્તિ~$s$ લો કે
        $s(x)=\Value{t}{M}$. \olref[ass]{prop:sat-quant}થી
        $\Sat{M}{!A(x)}[s]$. \olref[ext]{prop:ext-formulas}થી
        $\Sat{M}{!A(t)}[s]$. $!A(t)$ !!a{sentence} હોવાથી
        \olref[ass]{prop:sentence-sat-true}થી $\Sat{M}{!A(t)}$.}}{}
  \end{tagenumerate}
\end{proof}

\begin{probtag}{probEx,probAll}
  \olref[fol][syn][sem]{prop:quant-terms}ની સાબિતી પૂરી કરો.
\end{probtag}

\end{document}
